\documentclass[10pt,a4paper]{article}

% --- XeLaTeX: fonts + Russian ---
\usepackage{fontspec}
\usepackage{polyglossia}
\setdefaultlanguage{russian}
\setotherlanguage{english}
\setmainfont{PT Sans}
\setmonofont[Scale=MatchLowercase]{PT Mono}
\newfontfamily\cyrillicfonttt[Scale=MatchLowercase]{PT Mono}

% --- layout ---
\usepackage[a4paper,top=0.7cm,bottom=0.7cm,left=1.3cm,right=1.3cm]{geometry}
\usepackage{titlesec}
\usepackage{enumitem}
\usepackage{xcolor}
\usepackage[hidelinks]{hyperref}

\definecolor{rule}{HTML}{333333}

% section headers: bold, small rule under
\titleformat{\section}{\normalsize\bfseries}{}{0em}{\MakeUppercase}[{\color{rule}\titlerule[0.6pt]}]
\titlespacing*{\section}{0pt}{4pt}{2.2pt}

% tight bullet lists
\setlist[itemize]{leftmargin=1.15em, labelsep=0.5em, itemsep=1.2pt, topsep=2pt, parsep=0pt, before=\vspace{1pt}}
\renewcommand{\labelitemi}{\textbullet}

\setlength{\parindent}{0pt}
\setlength{\parskip}{0pt}
\pagestyle{empty}
\linespread{0.96}
\tolerance=1500
\emergencystretch=2.5em

% role line: bold title left, date right
\newcommand{\role}[2]{\textbf{#1}\hfill #2\par}

\begin{document}

% ===== Header =====
\begin{center}
  {\LARGE\bfseries Михайловский Илья Евгеньевич}\\[2pt]
  {\large Python Backend Developer}\\[3pt]
  Москва \textperiodcentered{} удалённо или гибрид \textperiodcentered{} английский C1, русский родной\\[2pt]
  Telegram: \href{https://t.me/Iluuuaaa}{@Iluuuaaa} \textperiodcentered{}
  \href{mailto:ilyamihailovsy@gmail.com}{ilyamihailovsy@gmail.com} \textperiodcentered{}
  \href{https://github.com/iluuua}{github.com/iluuua} \textperiodcentered{}
  +7 995 410-33-10
\end{center}

% ===== Summary =====
\section{О себе}
Разрабатываю backend на Python. Основал ReachFlow, спроектировал его серверную часть и довёл
до продакшена. Работаю с FastAPI, asyncpg, PostgreSQL, фоновыми процессами и асинхронными
интеграциями с Telegram и языковыми моделями. Сопровождаю развёртывание на Linux VPS.
Алгоритмическая подготовка: олимпиадные задачи на C++.

\section{Навыки}
\begin{itemize}
  \item Backend: Python, FastAPI, Pydantic, asyncio, REST API, JWT, OAuth
  \item Данные: PostgreSQL, SQL, asyncpg, pgvector, Redis
  \item Очереди и фоновые задачи: PostgreSQL, Redis и Dramatiq, отложенные действия и восстановление после сбоя
  \item Инфраструктура: Linux, Docker, Docker Compose, Caddy, nginx, systemd, GitHub Actions;
        базовое знакомство с Kubernetes (Deployment, CronJob, Secret)
  \item Качество и наблюдаемость: pytest, unittest, Ruff, проверки здоровья сервисов,
        структурированные логи, request-id и trace-id, OpenTelemetry
  \item Интеграции: Telegram API, MTProto, Telethon, Google OAuth, LLM через OpenRouter
  \item Алгоритмы: C++, алгоритмы и структуры данных
\end{itemize}

\section{Опыт работы}
\role{ReachFlow, Founder}{январь 2026 - настоящее время}
\textit{B2B SaaS для автоматизации продаж в Telegram \textperiodcentered{} \href{https://reachflow.tech}{reachflow.tech}}
Основатель продукта; отвечаю за backend, интеграции и эксплуатацию.
\begin{itemize}
  \item Спроектировал backend на FastAPI и asyncpg с изоляцией рабочих пространств,
        RLS-политиками PostgreSQL, JWT и Google OAuth. Подпись токенов проверяется по JWKS.
        Реализовал инициализацию базы, миграции и проверки схемы.
  \item Разработал фоновую обработку диалогов: группировку входящих сообщений, отложенные
        follow-up и восстановление задач после сбоя. Работал с Redis и Dramatiq;
        текущие API, шлюз и workers используют отдельные роли и очереди PostgreSQL.
        Интеграция с Telegram через Telethon включает вход по MTProto, хранение и восстановление
        сессий, MTProxy и SOCKS, лимиты отправки и обработку FloodWait.
  \item Разработал выполнение LLM-задач через OpenRouter: реестр 25 задач и 22 схем Pydantic
        с проверкой при импорте, профили моделей и резервные провайдеры без повторных попыток
        с одной и той же комбинацией модели и маршрута. Каждая попытка записывается
        с моделью, временем выполнения и стоимостью.
  \item Вынес чтение сайтов в сервис на FastAPI, Playwright и Chromium после замера пика памяти
        в 413 МБ при лимите 1 ГБ. Backend и рендерер независимо защищены от SSRF:
        проверяют DNS, публичность всех адресов и каждый редирект; разрешены порты 80 и 443.
  \item Настроил Docker Compose, Caddy/nginx, GitHub Actions с тестами и Ruff,
        проверки здоровья до уровня схемы базы, структурированные логи и OpenTelemetry.
        Работал с развёртыванием на Amvera и Linux VPS; для workers подготовлены примеры
        Kubernetes-манифестов.
  \item Сопровождаю автоматический деплой через GitHub webhooks и systemd.
        До миграций выполняются резервное копирование и восстановление на тестовой базе;
        после релиза проверяются версия образа и здоровье сервисов, при сбое предусмотрен откат.
\end{itemize}
\vspace{3pt}

\role{Delnovo, Chief Technology Officer}{июль 2023 - май 2026}
\begin{itemize}
  \item Разрабатывал коммерческих ботов для поддержки и продаж. Для дистрибьютора пластиковых окон
        реализовал асинхронного бота на Aiogram, интеграцию с Битрикс24 по REST API,
        RAG на PostgreSQL и уведомления о статусе заказа. По оценке заказчика бот автоматизировал
        около 80\% типовых консультаций.
  \item Разработал Telegram-бота для школы (\texttt{@shkolo\_teacher\_bot}): роли ученика,
        учителя и администратора, расписание, новости и отчётность. Настроил запуск в Docker.
\end{itemize}
\vspace{3pt}

\role{NEAR AI, Solutions Engineer}{июнь 2023 - июль 2024}
\begin{itemize}
  \item Решал алгоритмические задачи на C++ для обучения децентрализованного ИИ:
        оптимизировал время и память, оценивал асимптотику и улучшал структуру кода.
  \item Оформлял решения в LaTeX и проверял код участников международной команды.
\end{itemize}

\section{Проекты}
\role{Лидеры цифровой трансформации 2026, задача №5}{2026}
Капитан команды xixi-square. Разработал алгоритм обнаружения препятствий по облакам точек
3D-лидара Pandar128 и интеграцию с ROS 2 Humble. Отдельный процесс обрабатывает последний
кадр и пропускает устаревшие при перегрузке. Решение упаковано в Docker.
На синтетике организаторов найдены все 8 препятствий; на проверенных 15,7 км реальных записей
ложных тревог не было.

\vspace{3pt}
\textbf{Научная работа.} Написал статью по молекулярно-динамическому моделированию алюминиевой
матрицы с включением Fe$_4$Al$_{13}$. Код и материалы:
\href{https://github.com/iluuua/physics_md_al_fe}{physics\_md\_al\_fe}.

\section{Образование}
\textbf{Московский политехнический университет}, транспортный факультет, программа
<<Прикладная математика и информатика>>, выпуск 2029. \textbf{НИУ ВШЭ, ФКН},
вольнослушатель курсов той же программы.
\textbf{ОЦ <<Сириус>>}, <<Январская научная школа по математике и программированию>>,
92 академических часа, январь 2025. \textbf{Deep Learning School} ФПМИ МФТИ, 2025.
\textbf{Т-Образование}, <<Т-Поколение>>, курс <<Алгоритмы>>, параллель B, 2025.


\end{document}
